\documentclass[11pt,oneside]{article}

\usepackage[T1]{fontenc}
\usepackage[utf8]{inputenc}
\usepackage{lmodern}
\usepackage{microtype}
\usepackage[a4paper,margin=28mm,headheight=15pt]{geometry}
\usepackage{amsmath,amssymb,amsthm,mathtools}
\usepackage{booktabs,longtable,tabularx,array}
\usepackage{enumitem}
\usepackage{xcolor}
\usepackage{framed}
\usepackage{fancyhdr}
\usepackage{hyperref}
\usepackage{cleveref}
\usepackage{listings}
\usepackage{caption}

\definecolor{deepblue}{RGB}{18,60,105}
\definecolor{softblue}{RGB}{232,240,249}
\definecolor{softgray}{RGB}{245,246,247}
\definecolor{deepred}{RGB}{128,24,24}
\definecolor{deepgreen}{RGB}{18,96,58}

\hypersetup{
  colorlinks=true,
  linkcolor=deepblue,
  citecolor=deepblue,
  urlcolor=deepblue,
  pdftitle={The Two-Base Integer-Exponent Problem: Independent Audit and New Arithmetic Obstructions},
  pdfauthor={OpenAI, prepared for Vladimir Reshetnikov},
  pdfsubject={Alaoglu--Erdos integer-exponent conjecture},
  pdfkeywords={Alaoglu--Erdos, four exponentials, continued fractions, divided differences, integer jets}
}

\pagestyle{fancy}
\fancyhf{}
\lhead{The two-base integer-exponent problem}
\rhead{Independent audit}
\cfoot{\thepage}

\newtheorem{theorem}{Theorem}[section]
\newtheorem{proposition}[theorem]{Proposition}
\newtheorem{lemma}[theorem]{Lemma}
\newtheorem{corollary}[theorem]{Corollary}
\newtheorem{conjecture}[theorem]{Conjecture}
\theoremstyle{definition}
\newtheorem{definition}[theorem]{Definition}
\newtheorem{problem}[theorem]{Research target}
\theoremstyle{remark}
\newtheorem{remark}[theorem]{Remark}
\newtheorem{warning}[theorem]{Warning}

\newenvironment{claimbox}[1]{%
  \def\FrameCommand{\fcolorbox{deepblue}{softblue}}%
  \MakeFramed{\advance\hsize-2\width\FrameRestore}%
  \noindent\textbf{#1}\par\smallskip
}{\endMakeFramed}

\newenvironment{redbox}[1]{%
  \def\FrameCommand{\fcolorbox{deepred}{softgray}}%
  \MakeFramed{\advance\hsize-2\width\FrameRestore}%
  \noindent\textbf{#1}\par\smallskip
}{\endMakeFramed}

\newcommand{\Z}{\mathbb Z}
\newcommand{\N}{\mathbb N}
\newcommand{\Q}{\mathbb Q}
\newcommand{\R}{\mathbb R}
\newcommand{\C}{\mathbb C}
\newcommand{\ord}{\operatorname{ord}}
\newcommand{\rad}{\operatorname{rad}}
\newcommand{\sgn}{\operatorname{sgn}}
\newcommand{\lcm}{\operatorname{lcm}}
\newcommand{\vp}[2]{v_{#1}(#2)}
\newcommand{\AEconj}{Alaoglu--Erd\H{o}s}
\newcommand{\fourE}{Four Exponentials Conjecture}
\newcommand{\abs}[1]{\left|#1\right|}
\newcommand{\floor}[1]{\left\lfloor#1\right\rfloor}
\newcommand{\ceil}[1]{\left\lceil#1\right\rceil}
\newcommand{\e}{\mathrm e}
\newcommand{\dd}{\,\mathrm d}
\newcommand{\bigO}{\mathrm O}

\lstset{
  basicstyle=\ttfamily\small,
  frame=single,
  backgroundcolor=\color{softgray},
  breaklines=true,
  showstringspaces=false,
  columns=fullflexible
}

\title{\textbf{The Two-Base Integer-Exponent Problem}\\[0.3em]
\Large An Independent Proof Audit, Unimodular Gap Arithmetic,\\
Rational Secant Approximants, and a Fixed-Order Integer-Jet Barrier}
\author{Prepared for Vladimir Reshetnikov}
\date{August 21, 2026}

\begin{document}
\maketitle

\begin{abstract}
The \AEconj{} conjecture asks whether a real number $x$ satisfying
$2^x,3^x\in\Z$ must be an integer.  The attached source report was audited
line by line at the level of its mathematical claims.  It contains a large
body of exact reductions, formalized partial results, conditional structure,
and failed or incomplete research programs, but it explicitly does not
contain an unconditional proof.  The present document is deliberately
nonduplicative: it does not reproduce that corpus.  Instead it gives an
independent proof ledger and develops four additions.

First, a general unimodular collision theorem controls common divisors of
binomial gaps $B^q-A^p$.  Applied to consecutive continued-fraction
convergents to $\log 3/\log 2$, it proves that the gaps
$3^{q_j}-2^{p_j}$ are pairwise coprime at adjacent indices.  Second, under a
hypothetical counterexample, the corresponding output gaps produce positive
rational divided differences and rational secant approximants to $x$ that
alternate rigorously around $x$; their reduced denominator supports obey the
same unimodular separation.  Third, a finite order-signature theorem shows
that no finite collection of strict monomial comparisons can distinguish the
pair $(2,3)$ from infinitely many unrelated coprime integer pairs.  Fourth, a
matrix divisibility theorem proves that fixed-order Euler or Hasse integer
jets contribute only $O(RJ\log L)$ to the logarithmic arithmetic budget of an
$R\times R$ determinant with exponent scale $L$ and jet order $J$; therefore
fixed-order jets cannot alter an $R^2$ leading constant when $\log L=o(R)$.

These results are rigorous and new relative to the attached report, but they
do not close the final algebraic-independence gap.  No unconditional proof of
the \AEconj{} conjecture is claimed.  The exact missing implication is isolated,
and a complete conditional proof from the \fourE{} is included.
\end{abstract}

\clearpage
\begin{claimbox}{Outcome and claim ledger}
\noindent\textbf{Unconditional proof of the full conjecture.}\par
Not obtained; no such proof is present in the attached report.\medskip

\noindent\textbf{Conditional proof from Four Exponentials.}\par
Complete and elementary once the conjectural input is granted.\medskip

\noindent\textbf{Counterexample restrictions.}\par
Any counterexample is positive, irrational, transcendental, and gives multiplicatively independent integer outputs.\medskip

\noindent\textbf{New arithmetic theorem.}\par
Unimodular exponent pairs force a sharp gcd formula for binomial gaps.\medskip

\noindent\textbf{New approximation theorem.}\par
Hypothetical output gaps give rational secant approximants alternating around $x$, with separated denominator support.\medskip

\noindent\textbf{New no-go theorem.}\par
Finitely many monomial order comparisons cannot characterize the special logarithmic slope.\medskip

\noindent\textbf{New determinant barrier.}\par
Fixed-order entrywise integer jets are lower-order at the $R^2$ scale unless $J\log L\gtrsim R$ or a genuinely global cancellation mechanism is found.
\end{claimbox}

\tableofcontents
\newpage

\section{The problem and the standard boundary}

\subsection{Statement}

\begin{conjecture}[\AEconj{}]
\label{conj:AE}
If $x\in\R$ and $2^x,3^x\in\Z$, then $x\in\Z$.
\end{conjecture}

Because $2^x$ and $3^x$ are positive, the integer hypothesis means that they
are positive integers.  Set
\begin{equation}
  M:=2^x,\qquad N:=3^x,
  \label{eq:MN}
\end{equation}
whenever a hypothetical solution is under discussion, and put
\begin{equation}
  \alpha:=\frac{\log 3}{\log 2}=\log_2 3.
  \label{eq:alpha}
\end{equation}
Then $N=M^\alpha$, and the conjecture is equivalently the assertion that the
only positive integers $M$ for which $M^\alpha$ is an integer are powers of
$2$.

\subsection{Why a source audit matters}

The attached source file
\begin{center}
\texttt{/Article/alaoglu-erdos-unified-report/alaoglu-erdos-unified-report.tex}
\end{center}
is a very large unified research report.  Its title page and several later
status statements explicitly say that the full conjecture remains open.  It
contains many exact lemmas, conditional classifications, determinant
constructions, computational exclusions, and proposed interfaces, but the
existence of such a large body of correct partial work is not itself a proof.
The decisive audit question is whether any section derives a contradiction
from the sole assumptions in \cref{conj:AE}.  It does not.

The present report therefore follows two rules.
\begin{enumerate}[label=(\roman*)]
\item It does not copy or paraphrase at length the attached report's existing
  twenty-part survey, exact conditional monoid classification, determinant
  hierarchies, $p$-adic transport calculations, special-function audits, or
  formalization catalogue.
\item Every new claim below is proved from explicitly stated hypotheses, and
  every point at which a known transcendence conjecture would be needed is
  marked as such.
\end{enumerate}

\section{Elementary reductions that every proof must respect}

\begin{proposition}[Sign and small interval]
\label{prop:sign}
If $2^x\in\Z$, then either $x=0$ or $x\ge 1$.  In particular, a nonintegral
solution of \cref{conj:AE} would satisfy $x>1$.
\end{proposition}

\begin{proof}
If $x<0$, then $0<2^x<1$, so $2^x$ cannot be a positive integer.  If
$0<x<1$, then $1<2^x<2$, again impossible.  The remaining possibilities are
$x=0$ and $x\ge1$.  A nonintegral solution cannot equal $1$.
\end{proof}

\begin{proposition}[Rational exponents]
\label{prop:rational}
Let $a>1$ be an integer.  If $x\in\Q$ and $a^x\in\Z$, then the denominator of
$x$ in lowest terms divides every prime valuation of $a$.  In particular, if
$a=2$ and $2^x\in\Z$, then $x\in\Z$.
\end{proposition}

\begin{proof}
Write $x=r/s$ with $r\in\Z$, $s\in\N$, and $\gcd(r,s)=1$.  The case $r<0$ is
excluded when $a^x$ is a positive integer other than zero.  Put $a^x=m\in\N$.
Then $m^s=a^r$.  For every prime $p$,
\[
  s\,v_p(m)=r\,v_p(a).
\]
Since $\gcd(r,s)=1$, one has $s\mid v_p(a)$ for every $p\mid a$.  For $a=2$,
$v_2(a)=1$, hence $s=1$.
\end{proof}

\begin{corollary}
\label{cor:irrational}
Every nonintegral solution of \cref{conj:AE} is irrational.
\end{corollary}

\begin{proposition}[Algebraic exponents]
\label{prop:transcendental}
Every nonintegral solution of \cref{conj:AE} is transcendental.
\end{proposition}

\begin{proof}
By \cref{cor:irrational}, such an $x$ is irrational.  If it were algebraic,
then the Gelfond--Schneider theorem would make $2^x$ transcendental, contrary
to $2^x\in\Z$.
\end{proof}

\begin{proposition}[Output multiplicative independence]
\label{prop:independent}
Let $x>0$ and let $M,N$ be as in \eqref{eq:MN}.  Then $M$ and $N$ are
multiplicatively independent: if $M^rN^s=1$ for $r,s\in\Z$, then $r=s=0$.
\end{proposition}

\begin{proof}
The relation gives $2^{xr}3^{xs}=1$, hence $2^r3^s=1$ after taking the
positive $x$th root.  Unique factorization gives $r=s=0$.
\end{proof}

\subsection{Equivalent logarithmic formulations}

For positive integers $M,N$, the equality
\[
  \frac{\log M}{\log 2}=\frac{\log N}{\log 3}
\]
is equivalent to
\begin{equation}
  \det\begin{pmatrix}
  \log 2&\log M\\
  \log 3&\log N
  \end{pmatrix}=0.
  \label{eq:logdet}
\end{equation}
Thus the problem asks for a very special classification of rank-one
$2\times2$ matrices of logarithms of positive integers.  Baker's theorem
controls nonzero \emph{linear} forms in logarithms with algebraic
coefficients.  Equation \eqref{eq:logdet} is instead a bilinear relation among
four logarithms.  Treating the coefficients $\log M$ and $\log N$ as if they
were algebraic is precisely the invalid step in several tempting arguments.

\section{The exact Four Exponentials implication}

\begin{conjecture}[Four Exponentials]
\label{conj:4e}
Let $u_1,u_2$ be linearly independent over $\Q$, and let $v_1,v_2$ be
linearly independent over $\Q$.  Then at least one of the four numbers
$\exp(u_iv_j)$, $1\le i,j\le2$, is transcendental.
\end{conjecture}

\begin{theorem}[Conditional resolution]
\label{thm:4e}
The \fourE{} implies the \AEconj{} conjecture.
\end{theorem}

\begin{proof}
Assume $2^x$ and $3^x$ are integers.  If $x\notin\Q$, take
\[
  (u_1,u_2)=(\log2,\log3),\qquad (v_1,v_2)=(1,x).
\]
The first pair is $\Q$-linearly independent because $2$ and $3$ are
multiplicatively independent; the second pair is $\Q$-linearly independent
because $x$ is irrational.  Yet all four exponentials
\[
  \e^{\log2}=2,\quad \e^{\log3}=3,\quad
  \e^{x\log2}=2^x,\quad \e^{x\log3}=3^x
\]
are algebraic, contradicting \cref{conj:4e}.  Therefore $x\in\Q$, and
\cref{prop:rational} gives $x\in\Z$.
\end{proof}

\begin{redbox}{The unavoidable status conclusion}
The argument above is complete, but its decisive input is the unproved
\fourE{}.  Replacing that conjecture by a sentence such as ``the four
exponentials are algebraic, hence the rows or columns are dependent'' is not a
proof; it is exactly the conjectural implication.  Any unconditional solution
must either prove this special case of Four Exponentials or exploit the
integrality of the four exponentials in a way not available for arbitrary
algebraic values.
\end{redbox}

\section{Unimodular arithmetic of binomial gaps}
\label{sec:gaps}

The first new contribution is a general gcd theorem for gaps $B^q-A^p$.  It
is elementary, but it organizes a substantial amount of local arithmetic in
one determinant parameter.

\subsection{The collision determinant}

\begin{definition}
For exponent pairs $(p,q),(r,s)\in\N^2$, define their determinant distance by
\[
  \Delta((p,q),(r,s)):=ps-rq,
  \qquad D:=\abs{ps-rq}.
\]
The pairs are \emph{unimodular neighbors} when $D=1$.
\end{definition}

\begin{theorem}[Unimodular gap gcd theorem]
\label{thm:gapgcd}
Let $A,B\ge2$ be coprime integers, let $p,q,r,s\in\N$, and assume
$D=\abs{ps-rq}>0$.  Then
\begin{equation}
 \gcd(B^q-A^p,\,B^s-A^r)
 \mid \gcd(A^D-1,\,B^D-1).
 \label{eq:gapgcd}
\end{equation}
In particular, if $D=1$, then the sharp identity
\begin{equation}
 \gcd(B^q-A^p,\,B^s-A^r)
 = \gcd(A-1,B-1)
 \label{eq:unimodular}
\end{equation}
holds.
\end{theorem}

\begin{proof}
Let $d$ be the gcd on the left of \eqref{eq:gapgcd}.  Since $\gcd(A,B)=1$,
$d$ is coprime to $AB$: a prime dividing both $d$ and $A$ would also divide
$B^q$, and similarly with $A$ and $B$ interchanged.  Hence $A$ and $B$ are
units modulo $d$.  We have
\[
  B^q\equiv A^p\pmod d,\qquad B^s\equiv A^r\pmod d.
\]
Raise the first congruence to the power $s$ and the second to the power $q$.
The left sides agree, so
\[
 A^{ps}\equiv A^{rq}\pmod d,
\]
and therefore $A^{ps-rq}\equiv1\pmod d$.  The same argument, raising the
first congruence to $r$ and the second to $p$, gives
$B^{rq-sp}\equiv1\pmod d$.  Inverting when necessary yields
$A^D\equiv B^D\equiv1\pmod d$, proving \eqref{eq:gapgcd}.  If $D=1$,
this shows that the left-hand gcd divides $\gcd(A-1,B-1)$.  Conversely,
every divisor of both $A-1$ and $B-1$ divides each gap $B^q-A^p$ because
$A\equiv B\equiv1$ modulo that divisor.  This proves the equality
\eqref{eq:unimodular}.
\end{proof}

\begin{remark}[Primewise version without coprimality]
\label{rem:primewise}
If $A$ and $B$ are not assumed coprime, the same proof applies to every prime
$\ell\nmid AB$.  Thus the part of the left-hand gcd supported away from
$AB$ divides $\gcd(A^D-1,B^D-1)$.  Equivalently, any such common prime
satisfies
\[
  \ord_\ell(A)\mid D,\qquad \ord_\ell(B)\mid D.
\]
\end{remark}

\subsection{The continued-fraction recurrence shadow}

The determinant proof gives the cleanest gcd statement.  There is also an
exact recurrence explaining why adjacent gcds propagate unchanged along a
continued fraction.

\begin{proposition}[Gap recurrence]
\label{prop:gaprecurrence}
Suppose exponent pairs satisfy
\[
 p_{j+1}=a_{j+1}p_j+p_{j-1},\qquad
 q_{j+1}=a_{j+1}q_j+q_{j-1}
\]
with $a_{j+1}\in\N$.  Put $h_j=B^{q_j}-A^{p_j}$.  Then
\begin{align}
 h_{j+1}
 &=B^{q_{j-1}}h_j
   \sum_{t=0}^{a_{j+1}-1}
   B^{q_j(a_{j+1}-1-t)}A^{p_jt}
   +A^{a_{j+1}p_j}h_{j-1}.
 \label{eq:gaprecurrence}
\end{align}
If $\gcd(A,B)=1$, then
\begin{equation}
 \gcd(h_{j+1},h_j)=\gcd(h_j,h_{j-1}).
 \label{eq:gcdconservation}
\end{equation}
\end{proposition}

\begin{proof}
Substitute the recurrences for $p_{j+1},q_{j+1}$ and add and subtract
$B^{q_{j-1}}A^{a_{j+1}p_j}$:
\begin{align*}
 h_{j+1}
 &=B^{q_{j-1}}(B^{q_j})^{a_{j+1}}
   -A^{p_{j-1}}(A^{p_j})^{a_{j+1}}\\
 &=B^{q_{j-1}}
   \bigl((B^{q_j})^{a_{j+1}}-(A^{p_j})^{a_{j+1}}\bigr)
   +A^{a_{j+1}p_j}(B^{q_{j-1}}-A^{p_{j-1}}).
\end{align*}
Factor the difference of powers to obtain \eqref{eq:gaprecurrence}.  Modulo
$h_j$ this gives
\[
 h_{j+1}\equiv A^{a_{j+1}p_j}h_{j-1}\pmod{h_j}.
\]
When $\gcd(A,B)=1$, every $h_j=B^{q_j}-A^{p_j}$ is coprime to $A$.
Multiplication by $A^{a_{j+1}p_j}$ therefore does not change the gcd with
$h_j$, proving \eqref{eq:gcdconservation}.
\end{proof}

\begin{corollary}[Exact adjacent gcd along the $2$--$3$ slope]
\label{cor:exactadjacentgeneral}
Let $(p_j/q_j)$ be the convergents to $\log3/\log2$.  For every coprime
$A,B\ge2$ and every $j\ge0$,
\[
 \gcd(B^{q_j}-A^{p_j},\,B^{q_{j+1}}-A^{p_{j+1}})
 =\gcd(A-1,B-1).
\]
\end{corollary}

\begin{proof}
Consecutive convergents are unimodular, so this is already
\eqref{eq:unimodular}.  Alternatively, \cref{prop:gaprecurrence} shows the
adjacent gcd is constant; at the first two convergents $(1,1)$ and $(2,1)$ it
is
\[
 \gcd(B-A,B-A^2)=\gcd(A-1,B-1),
\]
where coprimality removes the possible factor from $A$.
\end{proof}

\subsection{Continued fractions and adjacent coprimality}

Let
\[
  \frac{p_j}{q_j}=[a_0;a_1,\ldots,a_j]
\]
be the convergents to $\alpha=\log3/\log2$, indexed from $j=0$.  Consecutive
convergents satisfy
\begin{equation}
  \abs{p_jq_{j+1}-p_{j+1}q_j}=1.
  \label{eq:cfdet}
\end{equation}
Define the signed base gaps
\begin{equation}
  g_j:=3^{q_j}-2^{p_j}.
  \label{eq:basegap}
\end{equation}

\begin{theorem}[Adjacent base gaps are coprime]
\label{thm:adjacent}
For every $j\ge0$,
\begin{equation}
  \gcd(g_j,g_{j+1})=1.
  \label{eq:adjacentcoprime}
\end{equation}
\end{theorem}

\begin{proof}
Apply \cref{thm:gapgcd} with $(A,B)=(2,3)$ and the two consecutive convergent
pairs.  By \eqref{eq:cfdet}, $D=1$, so the gcd divides
$\gcd(2-1,3-1)=1$.
\end{proof}

\begin{corollary}[Higher-distance collision spectrum]
\label{cor:higherdistance}
For any $j\ne k$, put
\[
  D_{jk}:=\abs{p_jq_k-p_kq_j}.
\]
Then
\[
  \gcd(g_j,g_k)\mid\gcd(2^{D_{jk}}-1,3^{D_{jk}}-1).
\]
Thus every repeated prime divisor of two base gaps has multiplicative orders
of both $2$ and $3$ dividing the determinant distance $D_{jk}$.
\end{corollary}

The first convergents and gaps are shown in \cref{tab:gaps}.  The huge growth
of the gaps is compatible with their adjacent coprimality; relative
Diophantine closeness does not mean small absolute difference once the
exponential scale is restored.

\begin{table}[ht]
\centering
\caption{Initial convergents to $\log3/\log2$ and signed gaps.}
\label{tab:gaps}
\begin{tabular}{@{}rrrrl@{}}
\toprule
$j$ & $p_j$ & $q_j$ & $\sgn(g_j)$ & $\abs{g_j}$ \\
\midrule
0 & 1 & 1 & $+$ & $1$\\
1 & 2 & 1 & $-$ & $1$\\
2 & 3 & 2 & $+$ & $1$\\
3 & 8 & 5 & $-$ & $13$\\
4 & 19 & 12 & $+$ & $7153$\\
5 & 65 & 41 & $-$ & $420491770248316829$\\
6 & 84 & 53 & $+$ & $40432553845953101497907$\\
7 & 485 & 306 & $-$ & a $144$-digit integer\\
8 & 1054 & 665 & $+$ & a $313$-digit integer\\
\bottomrule
\end{tabular}
\end{table}

\subsection{Output gaps under a hypothetical counterexample}

Assume now that $x>1$ is a nonintegral solution, and retain $M=2^x$,
$N=3^x$.  Define
\begin{equation}
  G_j:=N^{q_j}-M^{p_j}.
  \label{eq:outputgap}
\end{equation}
The common logarithmic slope implies
\begin{equation}
  \sgn(G_j)=\sgn(g_j).
  \label{eq:samesign}
\end{equation}
Indeed, if
\begin{equation}
  \delta_j:=q_j\log3-p_j\log2,
  \label{eq:delta}
\end{equation}
then
\[
  g_j=2^{p_j}(\e^{\delta_j}-1),\qquad
  G_j=M^{p_j}(\e^{x\delta_j}-1).
\]

\begin{corollary}[Output collision restriction]
\label{cor:outputgcd}
For consecutive convergents, if $\gcd(M,N)=1$, then the exact formula
\[
  \gcd(G_j,G_{j+1})=\gcd(M-1,N-1)
\]
holds.
Without the coprimality assumption, the part of
$\gcd(G_j,G_{j+1})$ supported away from $MN$ satisfies the same divisibility.
More generally, at indices $j\ne k$ every common prime $\ell\nmid MN$ has
$\ord_\ell(M),\ord_\ell(N)\mid D_{jk}$.
\end{corollary}

\begin{proof}
Apply \cref{thm:gapgcd} or \cref{rem:primewise} to $(A,B)=(M,N)$.
\end{proof}

\begin{remark}[What this does not prove]
The base gaps have adjacent gcd exactly $1$.  When $M,N$ are coprime, the
output gaps have the equally rigid but not contradictory constant adjacent
gcd $\gcd(M-1,N-1)$; without coprimality, additional factors supported on
$MN$ may occur.  There is no known theorem forcing the same prime to recur in adjacent
output gaps.  Consequently the gcd theorem supplies arithmetic separation,
not a contradiction.
\end{remark}

\section{Transported divided differences}
\label{sec:secants}

The same convergents produce a rational family whose analytic behavior is
completely rigid.

\subsection{Positive rational divided differences}

For any positive $p,q$ with $3^q\ne2^p$, define
\begin{equation}
  Q_{p,q}:=\frac{N^q-M^p}{3^q-2^p}.
  \label{eq:Qpq}
\end{equation}

\begin{theorem}[Exact divided-difference transport]
\label{thm:divdiff}
Under a hypothetical solution $x>1$, $Q_{p,q}$ is a positive rational number
and
\begin{equation}
 Q_{p,q}
 =x\int_0^1\bigl((1-t)2^p+t3^q\bigr)^{x-1}\dd t.
 \label{eq:integralQ}
\end{equation}
If $Q_{p,q}=a_{p,q}/b_{p,q}$ in lowest terms with $b_{p,q}>0$, then
\begin{equation}
 b_{p,q}=\frac{\abs{3^q-2^p}}
 {\gcd(\abs{3^q-2^p},\abs{N^q-M^p})},
 \qquad
 b_{p,q}\mid\abs{3^q-2^p}.
 \label{eq:denexact}
\end{equation}
\end{theorem}

\begin{proof}
Both numerator and denominator in \eqref{eq:Qpq} are nonzero integers.  Their
signs agree because $t\mapsto t^x$ is strictly increasing on the positive
reals, so the quotient is positive.  The identity
\[
  \frac{b^x-a^x}{b-a}=x\int_0^1((1-t)a+tb)^{x-1}\dd t
\]
follows from the fundamental theorem of calculus after the affine
substitution $u=(1-t)a+tb$; take $a=2^p$, $b=3^q$.  Formula
\eqref{eq:denexact} is the standard reduction of a ratio of two integers.
\end{proof}

For convergents, write $Q_j:=Q_{p_j,q_j}$ and let $b_j$ be its reduced
denominator.

\begin{corollary}[Adjacent denominator separation]
\label{cor:denomsep}
For every $j\ge0$,
\[
  \gcd(b_j,b_{j+1})=1.
\]
\end{corollary}

\begin{proof}
By \cref{thm:divdiff}, $b_j\mid\abs{g_j}$.  Apply
\cref{thm:adjacent}.
\end{proof}

\subsection{Asymptotics at convergent nodes}

The continued-fraction estimate gives
\begin{equation}
  \abs{\delta_j}
  =\log2\,\abs{q_j\alpha-p_j}
  <\frac{\log2}{q_{j+1}}.
  \label{eq:deltabound}
\end{equation}
Moreover, the signs of $\delta_j$ alternate.

\begin{proposition}[Asymptotic secant scale]
\label{prop:Qasymp}
As $j\to\infty$,
\begin{align}
 Q_j
 &=2^{p_j(x-1)}\frac{\e^{x\delta_j}-1}{\e^{\delta_j}-1}
 \label{eq:Qexactexp}\\
 &=x\,2^{p_j(x-1)}
 \left(1+\frac{x-1}{2}\delta_j+\bigO_x(\delta_j^2)\right)
 \label{eq:Qasymp}\\
 \log Q_j
 &=(x-1)p_j\log2+\log x+\bigO_x(q_{j+1}^{-1})
 \label{eq:logQ}
\end{align}
\end{proposition}

\begin{proof}
Use $3^{q_j}=2^{p_j}\e^{\delta_j}$ and
$N^{q_j}=M^{p_j}\e^{x\delta_j}$ in \eqref{eq:Qpq}.  Since
$M^{p_j}/2^{p_j}=2^{p_j(x-1)}$, this gives \eqref{eq:Qexactexp}.  Taylor
expansion at $\delta=0$ yields
\[
 \frac{\e^{x\delta}-1}{\e^\delta-1}
 =x+\frac{x(x-1)}2\delta
 +\frac{x(x-1)(2x-1)}{12}\delta^2+\bigO_x(\delta^3).
\]
Combine this with \eqref{eq:deltabound}; taking logarithms is legitimate for
all sufficiently large $j$ because the parenthetical factor tends to $1$.
\end{proof}

\subsection{Rational approximants that bracket the unknown exponent}

Normalize away the exponential scale by defining
\begin{equation}
 T_j:=\frac{2^{p_j}}{M^{p_j}}Q_j
 =\frac{2^{p_j}(N^{q_j}-M^{p_j})}
 {M^{p_j}(3^{q_j}-2^{p_j})}.
 \label{eq:Tj}
\end{equation}
Every $T_j$ is rational.  Define the analytic secant function
\begin{equation}
 \Phi_x(\delta):=
 \begin{cases}
 \dfrac{\e^{x\delta}-1}{\e^\delta-1},&\delta\ne0,\\[0.8em]
 x,&\delta=0.
 \end{cases}
 \label{eq:Phi}
\end{equation}
Then $T_j=\Phi_x(\delta_j)$.

\begin{theorem}[Alternating rational secant brackets]
\label{thm:brackets}
Let $x>1$.  The function $\Phi_x$ is strictly increasing on $\R$.  Hence
\begin{equation}
  \sgn(T_j-x)=\sgn(\delta_j),
  \label{eq:bracketsign}
\end{equation}
and consecutive $T_j$ lie on opposite sides of $x$.  Quantitatively,
\begin{equation}
  T_j-x=\frac{x(x-1)}2\delta_j+\bigO_x(\delta_j^2)
  =\bigO_x(q_{j+1}^{-1}).
  \label{eq:Terror}
\end{equation}
Thus a hypothetical counterexample generates an explicit alternating sequence
of rational upper and lower bounds converging to $x$.
\end{theorem}

\begin{proof}
For every real $\delta$,
\begin{equation}
 \Phi_x(\delta)
 =x\int_0^1(1-t+t\e^\delta)^{x-1}\dd t.
 \label{eq:Phiintegral}
\end{equation}
Differentiating under the integral sign gives
\[
 \Phi_x'(\delta)
 =x(x-1)\e^\delta
 \int_0^1t(1-t+t\e^\delta)^{x-2}\dd t>0.
\]
Therefore $\Phi_x(\delta)-\Phi_x(0)$ has the sign of $\delta$, proving
\eqref{eq:bracketsign}.  The expansion in \eqref{eq:Terror} is the Taylor
series already used in \cref{prop:Qasymp}, and the final estimate follows from
\eqref{eq:deltabound}.
\end{proof}

\begin{definition}[Denominator away from a fixed integer]
For a positive integer $d$ and $M\ge2$, let $d^{(M)}$ denote the largest
divisor of $d$ coprime to $M$.
\end{definition}

\begin{corollary}[External denominator separation for the rational brackets]
\label{cor:Textden}
Let $d_j$ be the reduced denominator of $T_j$.  Then
\[
 d_j^{(M)}\mid\abs{g_j},\qquad
 \gcd(d_j^{(M)},d_{j+1}^{(M)})=1.
\]
\end{corollary}

\begin{proof}
Equation \eqref{eq:Tj} is obtained from $Q_j$ by multiplying by
$2^{p_j}/M^{p_j}$.  New denominator primes can therefore come only from $M$;
away from $M$, the reduced denominator still divides the denominator $b_j$ of
$Q_j$, hence divides $g_j$.  Now use \cref{thm:adjacent}.
\end{proof}

\subsection{Why the rational brackets are not yet an irrationality-measure
contradiction}

The approximation error in \eqref{eq:Terror} is of order $q_{j+1}^{-1}$,
whereas the denominator of $T_j$ can contain $M^{p_j}$ and has exponential
height in $p_j$.  In terms of that denominator, the approximation is far too
weak to contradict any general irrationality measure.  More importantly,
$x$ is already forced to be transcendental by \cref{prop:transcendental}, so
algebraic irrationality measures do not apply.  The new information is not
the existence of rational approximants, but their canonical construction,
alternation, and prime-support separation.

\section{A finite order-signature no-go theorem}
\label{sec:signature}

The common logarithmic slope means that every monomial comparison is
transported:
\[
  2^a3^b<2^c3^d
  \quad\Longleftrightarrow\quad
  M^aN^b<M^cN^d.
\]
It is tempting to search for a finite list of such inequalities that only the
pair $(2,3)$ can satisfy.  The next theorem rules out that strategy in its
pure form.

\begin{theorem}[Finite order-signature indistinguishability]
\label{thm:signature}
Let $F\subset\Z^2\setminus\{(0,0)\}$ be finite.  There exist infinitely many
coprime, multiplicatively independent integer pairs $(M,N)$ with $M,N>1$ such
that
\begin{equation}
 \sgn(a\log M+b\log N)
 =\sgn(a\log2+b\log3)
 \quad\text{for every }(a,b)\in F.
 \label{eq:signature}
\end{equation}
\end{theorem}

\begin{proof}
Write $\rho=\log N/\log M$.  Since
$\alpha=\log3/\log2$ is irrational, $a+b\alpha\ne0$ for every nonzero
$(a,b)\in\Z^2$.  Because $F$ is finite, there is an open interval $I$
containing $\alpha$ on which every function $\rho\mapsto a+b\rho$ has the
same sign as at $\alpha$.

Choose an arbitrarily large prime $P$ and set $M=P$.  Among any two
consecutive integers nearest to $P^\alpha$, at least one is not divisible by
$P$; choose such an integer as $N$.  Then $\gcd(M,N)=1$, and
\[
 \frac{N}{P^\alpha}\longrightarrow1,
 \qquad
 \frac{\log N}{\log P}\longrightarrow\alpha
 \quad(P\to\infty).
\]
For all sufficiently large choices, $\rho\in I$, so
\[
 \sgn(a\log M+b\log N)
 =\sgn(\log M)\,\sgn(a+b\rho)
 =\sgn(a+b\alpha),
\]
which is the right side of \eqref{eq:signature}.  The pair is
multiplicatively independent because two coprime integers greater than $1$
cannot satisfy a nontrivial integer power relation.  Infinitely many primes
$P$ give infinitely many pairs.
\end{proof}

\begin{corollary}[Finite monomial order patterns]
\label{cor:finiteorder}
For every finite set $E\subset\Z^2$, there exist infinitely many coprime,
multiplicatively independent pairs $(M,N)$ for which the entire strict order
pattern of the monomials $M^aN^b$, $(a,b)\in E$, is identical to that of
$2^a3^b$.
\end{corollary}

\begin{proof}
Apply \cref{thm:signature} to all nonzero differences of pairs in $E$.
\end{proof}

\begin{warning}
This theorem does not rule out a finite proof using congruences, exact
integrality, valuations, or transcendence input.  It rules out only a proof
whose distinguishing information is a finite collection of strict monomial
order comparisons.  Any order-theoretic approach must use an infinite or
uniform property, or combine order with genuinely arithmetic data.
\end{warning}

\section{A fixed-order integer-jet barrier}
\label{sec:jets}

Integer jets are attractive in determinant methods because Euler derivatives
preserve integrality at integer points:
\[
 (X\partial_X)^r(Y\partial_Y)^s(X^nY^k)
 =n^rk^sX^nY^k.
\]
The attached report identifies nonvanishing and scaling as unresolved.  The
next result gives a sharp general budget for the automatic divisibility that
can come from the entrywise integer multipliers alone.

\subsection{A matrix theorem}

For an integer matrix $C=(c_{ij})_{1\le i,j\le R}$, allow zero entries and
set $v_p(0)=+\infty$.  Let
\[
 \Sigma(C):=\{\sigma\in S_R:c_{i,\sigma(i)}\ne0\text{ for all }i\}.
\]
If $\Sigma(C)\ne\varnothing$, define
\begin{equation}
 \tau_p(C):=\min_{\sigma\in\Sigma(C)}
 \sum_{i=1}^R v_p(c_{i,\sigma(i)}),
 \qquad
 T(C):=\prod_p p^{\tau_p(C)}.
 \label{eq:tau}
\end{equation}
Only finitely many factors in $T(C)$ are nontrivial.

\begin{theorem}[Entrywise decoration budget]
\label{thm:decoration}
Let $E=(e_{ij})$ and $C=(c_{ij})$ be integer $R\times R$ matrices.  Assume
$\Sigma(C)\ne\varnothing$ and $\abs{c_{ij}}\le B$ for every nonzero entry,
where $B\ge1$.  Then
\begin{equation}
  T(C)\mid\det(E\circ C),
  \label{eq:Tdiv}
\end{equation}
where $\circ$ denotes entrywise product, and
\begin{equation}
  \log T(C)\le R\log B.
  \label{eq:Tbound}
\end{equation}
If additionally $\abs{e_{ij}}\le E_0$, then
\begin{equation}
 \log\abs{\det(E\circ C)}
 \le\log(R!)+R\log E_0+R\log B
 \label{eq:archbudget}
\end{equation}
whenever the determinant is nonzero.
\end{theorem}

\begin{proof}
Expand the determinant by permutations.  A term with a zero $c$-entry
vanishes.  Every nonzero term indexed by $\sigma\in\Sigma(C)$ has $p$-adic
valuation at least
\[
 \sum_i v_p(c_{i,\sigma(i)})\ge\tau_p(C),
\]
because the $e_{ij}$ are integers.  Hence every term, and therefore their sum,
is divisible by $p^{\tau_p(C)}$ for every $p$, proving \eqref{eq:Tdiv}.

For the size bound, use ``sum of minima is at most the minimum of sums'':
\begin{align*}
 \log T(C)
 &=\sum_p\tau_p(C)\log p\\
 &\le\min_{\sigma\in\Sigma(C)}
   \sum_p\sum_i v_p(c_{i,\sigma(i)})\log p\\
 &=\min_{\sigma\in\Sigma(C)}
   \sum_i\log\abs{c_{i,\sigma(i)}}
 \le R\log B.
\end{align*}
Finally, every Leibniz term of $\det(E\circ C)$ has absolute value at most
$(E_0B)^R$, and there are $R!$ terms, giving \eqref{eq:archbudget}.
\end{proof}

\begin{remark}[Optimality of the scale]
The bound $R\log B$ cannot be improved in general: take $C=B I_R$ and an
appropriate diagonal $E$.  The theorem is therefore a structural ceiling,
not merely a loose use of Hadamard's inequality.
\end{remark}

\subsection{Euler and Hasse jets}

Suppose exponents $n,k$ lie in $[0,L]\cap\Z$ and all Euler jet orders obey
$r+s\le J$.  The multiplier satisfies
\[
 \abs{n^rk^s}\le L^{r+s}\le L^J
\]
when nonzero.

\begin{corollary}[Euler-jet budget]
\label{cor:eulerbudget}
For an $R\times R$ determinant obtained by entrywise decorating an integral
evaluation matrix with Euler-jet multipliers of total order at most $J$ and
exponent scale at most $L$, the automatic extra divisor supplied solely by
those multipliers has logarithm at most
\begin{equation}
  RJ\log L.
  \label{eq:eulerbudget}
\end{equation}
The same quantity is the maximum added term in the crude logarithmic
archimedean bound.
\end{corollary}

\begin{proof}
Apply \cref{thm:decoration} with $B=L^J$.
\end{proof}

Hasse derivatives produce multipliers $\binom nr\binom ks$, also bounded by
$L^{r+s}\le L^J$.  They additionally lower the powers of $X$ and $Y$; those
row or column factors must be charged separately and cannot be counted as a
free jet gain.

\begin{corollary}[Leading-order no-go for fixed-order jets]
\label{cor:jetnogo}
Consider a sequence of $R\times R$ determinant constructions with exponent
scale $L_R$ and jet order $J_R$.  If
\begin{equation}
  J_R\log L_R=o(R),
  \label{eq:subcriticaljets}
\end{equation}
then the termwise guaranteed integer-jet divisor and the corresponding crude\
archimedean multiplier budget are both $o(R^2)$ in logarithmic scale.  In particular, fixed $J_R$ with polynomially growing $L_R$ cannot
alter an $R^2$ leading constant in an archimedean-versus-nonarchimedean
determinant comparison.
\end{corollary}

\begin{proof}
By \eqref{eq:eulerbudget}, the contribution is at most
$R J_R\log L_R=o(R^2)$.
\end{proof}

\begin{claimbox}{What a successful jet method must add}
At least one of the following must occur:
\begin{enumerate}[label=(\alph*)]
\item $J_R\log L_R\gtrsim R$, so that the jet order or exponent scale is
  genuinely large;
\item the determinant has systematic cancellation beyond the minimum
  valuation of its individual Leibniz terms;
\item confluence creates a global factorization or rank defect not representable
  as bounded entrywise integer decoration;
\item a nonlocal arithmetic identity couples several primes or several jet
  levels at once.
\end{enumerate}
Fixed-order differentiation alone cannot supply the missing leading-order
arithmetic constant.
\end{claimbox}

\section{Stress tests of possible closure arguments}
\label{sec:stress}

This section records the last steps one might hope to extract from the new
lemmas and explains precisely why each step is currently unavailable.

\subsection{Can adjacent coprimality force a contradiction?}

No.  The integers $g_j$ are allowed to introduce new prime divisors at every
index.  The theorem $\gcd(g_j,g_{j+1})=1$ is entirely compatible with their
exponential growth.  The reduced denominators $b_j$ of $Q_j$ inherit this
separation, but a sequence of rationals can have pairwise coprime adjacent
denominators indefinitely.  A contradiction would require a separate theorem
forcing denominator recurrence or confinement to a fixed finite prime set.
No such theorem follows from the hypotheses.

\subsection{Can one transfer primitive primes from base gaps to output gaps?}

The equality of logarithmic slopes synchronizes the signs and relative sizes
of $g_j$ and $G_j$, but it does not synchronize their residues modulo a prime.
A prime dividing $g_j$ satisfies $3^{q_j}\equiv2^{p_j}$, whereas a prime
dividing $G_j$ satisfies $N^{q_j}\equiv M^{p_j}$.  The real-logarithmic
relation gives no direct congruence between these two statements.  A theorem
forcing infinitely many common prime divisors between corresponding base and
output gaps would be a genuinely new local-global principle and would go
well beyond the elementary gcd calculations.

\subsection{Can the rational brackets prove \texorpdfstring{$x$}{x} rational?}

The approximants $T_j$ are canonical and alternate, but their error is only
$O(q_{j+1}^{-1})$ while their possible denominators are exponentially large.
They do not give a Liouville-type approximation.  Nor can Roth's theorem or
an algebraic irrationality measure be applied, because a counterexample is
already known to be transcendental.

\subsection{Can Pad\'e approximation close the logarithm-product relation?}

A Pad\'e construction becomes decisive only if its denominator clearing and
archimedean smallness beat each other at the same leading scale.  The relation
\eqref{eq:logdet} involves products of logarithms, not a nonzero linear form
with algebraic coefficients.  Recent Pad\'e advances for polylogarithmic
periods provide useful techniques, but they do not currently imply the
required special algebraic independence of
$\log2,\log3,\log M,\log N$.

\subsection{Can automaticity or Cobham-type rigidity be invoked?}

Cobham theorems require an actual finite automaton, recognizability in two
multiplicatively independent numeration systems, or an equivalent finite-state
hypothesis.  The integrality of $2^x$ and $3^x$ does not by itself make the
digit sequence of any associated quantity automatic.  Inferring automaticity
from the existence of two integer values is an unsupported bridge.

\subsection{Can finite order geometry identify the slope?}

No finite list of strict monomial comparisons can do so, by
\cref{thm:signature}.  An order-based proof must use a uniform statement over
unbounded exponents, a compactness principle retaining arithmetic data, or a
separate exact identity.

\subsection{Can fixed-order confluence repair a determinant constant?}

Not through the entrywise multiplier divisibility alone.  The bound in
\cref{cor:jetnogo} shows that such a gain is lower-order.  A successful
confluent determinant would need global cancellation, a growing jet regime,
or an exact factorization invisible in termwise valuations.

\section{Conditional and restricted positive results}

The full conjecture remains beyond current unconditional methods, but several
nearby statements are complete.

\begin{theorem}[Three algebraic bases from Six Exponentials]
\label{thm:sixexp}
Let $a_1,a_2,a_3>1$ be algebraic numbers whose logarithms are linearly
independent over $\Q$.  If $x\in\R$ and all three $a_i^x$ are algebraic, then
$x\in\Q$.
\end{theorem}

\begin{proof}
If $x\notin\Q$, then $1,x$ are $\Q$-linearly independent.  The Six
Exponentials Theorem applied to the three logarithms $\log a_i$ and the two
numbers $1,x$ says that at least one of the six exponentials
$a_i$ or $a_i^x$ is transcendental, contrary to the hypotheses.
\end{proof}

\begin{corollary}
If $2^x,3^x,5^x\in\Z$, then $x\in\Z$.
\end{corollary}

\begin{proof}
The logarithms of $2,3,5$ are $\Q$-linearly independent by unique
factorization.  Apply \cref{thm:sixexp}, then \cref{prop:rational}.
\end{proof}

\begin{proposition}[A logarithm-product independence criterion]
\label{prop:productcriterion}
Suppose the products $\log p\,\log q$, with $p,q$ ranging over primes, are
linearly independent over $\Q$ subject only to the symmetry
$\log p\log q=\log q\log p$.  Then \cref{conj:AE} holds.
\end{proposition}

\begin{proof}
Factor hypothetical integer outputs as
$M=\prod_p p^{a_p}$ and $N=\prod_q q^{b_q}$.  Equation \eqref{eq:logdet}
becomes
\[
 \sum_p a_p\log p\log3-
 \sum_q b_q\log q\log2=0.
\]
Under the stated independence, the coefficient of every unordered prime pair
must vanish.  The only possible supports are then $M=2^m$ and $N=3^m$ with
the same nonnegative integer $m$, giving $x=m$.
\end{proof}

\begin{remark}
The hypothesis in \cref{prop:productcriterion} is far stronger than any
currently known theorem and is stated only to identify the exact type of
algebraic independence that would settle the problem.
\end{remark}

\section{Precise research targets exposed by the new work}

\begin{problem}[Cross-family prime transfer]
Find a theorem that uses the exact equality
$\log N/\log M=\log3/\log2$ to force a nontrivial relation between the prime
supports of
\[
  3^{q_j}-2^{p_j}
  \quad\text{and}\quad
  N^{q_j}-M^{p_j}
\]
for infinitely many convergents.  Real asymptotics and within-family gcd
bounds are insufficient; the needed input must be genuinely local-global.
\end{problem}

\begin{problem}[Denominator confinement]
Prove, from the counterexample hypotheses, that the reduced denominators of
$Q_j$ or the prime-to-$M$ parts of the denominators of $T_j$ lie in a fixed
finite prime set or satisfy a recurrent divisibility law.  Combined with the
unimodular separation in \cref{cor:denomsep,cor:Textden}, such a result would
create strong pressure toward eventual integrality.
\end{problem}

\begin{problem}[Leading-order jet cancellation]
Construct a confluent determinant whose $p$-adic valuation exceeds the
minimum valuation of its nonzero Leibniz terms by $cR^2+o(R^2)$ for some
$c>0$, while preserving nonvanishing and a matching archimedean estimate.
Theorem~\ref{thm:decoration} shows that bounded entrywise integer multipliers
alone cannot do this.
\end{problem}

\begin{problem}[Infinite order rigidity]
Replace finite monomial comparisons by a uniform order property over an
unbounded region and prove that the combination of that property with
integrality characterizes the slope $\log3/\log2$.  Theorem~\ref{thm:signature}
shows why any finite truncation is inadequate.
\end{problem}

\section{Literature check through August 21, 2026}
\label{sec:literature}

The current literature check supports the same status conclusion as the
source audit.

\begin{enumerate}[label=(\roman*)]
\item Waldschmidt's chapter on the Four Exponentials Problem states the exact
  question whether $2^t$ and $3^t$ integral force $t$ to be a nonnegative
  integer, and explains that a positive answer follows from Four
  Exponentials; the relevant complex and $p$-adic questions remain open.
\item The 2026 Lean formalization of Gelfond--Schneider rigorously covers the
  algebraic-irrational case used in \cref{prop:transcendental}.  It does not
  cover transcendental exponents, which are exactly the remaining case.
\item Recent Pad\'e work on linear independence of polylogarithmic periods and
  products of polylogarithms develops powerful approximation machinery, but
  does not establish the logarithm-product independence required by
  \eqref{eq:logdet}.
\item Recent work on additive relations in irrational powers concerns a
  different rigidity problem.  It provides useful surrounding techniques but
  does not imply simultaneous integrality of the two powers only at integer
  exponents.
\end{enumerate}

This report makes no global priority claim for its elementary lemmas.  The
novelty statement is limited and verifiable: the particular unimodular gap gcd package, rational secant bracket package,\
finite order-signature theorem, and fixed-order entrywise jet budget were not\
located in the attached source report in the forms proved here.

\section{Conclusion}

The requested implication has not been proved unconditionally here, and it
would be mathematically false to present the partial results as a proof.  The
source audit and the current literature both place the problem at the Four
Exponentials boundary.  Nevertheless, the investigation produces four exact
advances relative to the attached report:
\begin{enumerate}[label=(\arabic*)]
\item determinant distance controls common divisors of exponential binomial
  gaps, giving adjacent coprimality along the continued fraction of
  $\log3/\log2$;
\item a hypothetical counterexample transports these gaps into rational
  divided differences and alternating rational secant brackets for $x$, with
  sharply separated denominator support;
\item finite monomial order data cannot isolate the special logarithmic
  slope;
\item fixed-order integer jets are provably lower-order at the termwise budget level\
  unless a new global cancellation mechanism is found.
\end{enumerate}

The first two results create a concrete arithmetic interface between
continued fractions and the hypothetical integer outputs.  The last two
eliminate broad classes of finite or local closing arguments.  The remaining
challenge is correspondingly narrow: one needs a mechanism coupling the real
logarithmic identity to prime support or determinant cancellation at a scale
that current linear-forms, Pad\'e, automatic, and termwise $p$-adic methods do
not reach.

\appendix

\section{Detailed asymptotic bounds}

For completeness, here is a nonasymptotic version of the estimates used in
\cref{sec:secants}.  If $\abs{u}\le1/2$, the mean value theorem gives
\[
  \e^{-1/2}\abs{u}\le\abs{\e^u-1}\le\e^{1/2}\abs{u}.
\]
Consequently, whenever $\abs{\delta_j}\le1/(2x)$,
\[
 x\e^{-1}\,2^{p_j(x-1)}
 \le Q_j \le
 x\e\,2^{p_j(x-1)},
\]
a deliberately coarse but uniform two-sided exponential estimate.  More
precise constants follow by applying Taylor's theorem with remainder to
$\Phi_x$ on a compact interval around zero.

The continued-fraction error bound
\[
 \abs{\alpha-\frac{p_j}{q_j}}
 <\frac{1}{q_jq_{j+1}}
\]
immediately gives \eqref{eq:deltabound}.  Since the convergents alternate
around an irrational number, $\delta_j$ alternates in sign, and the strict
monotonicity of $\Phi_x$ converts this into the rational brackets of
\cref{thm:brackets} without any asymptotic qualification.

\section{Reproducible continued-fraction computation}

The numerical entries in \cref{tab:gaps} can be reproduced with arbitrary
precision using the following short script.

\begin{lstlisting}[language=Python]
import mpmath as mp
from math import gcd

mp.mp.dps = 200
alpha = mp.log(3) / mp.log(2)

# Continued-fraction coefficients.
a = []
y = alpha
for _ in range(9):
    z = int(mp.floor(y))
    a.append(z)
    y = 1 / (y - z)

# Convergents p_j / q_j.
p0, p1 = 0, 1
q0, q1 = 1, 0
conv = []
for z in a:
    p = z * p1 + p0
    q = z * q1 + q0
    conv.append((p, q))
    p0, p1 = p1, p
    q0, q1 = q1, q

gaps = [pow(3, q) - pow(2, p) for p, q in conv]
for j, ((p, q), g) in enumerate(zip(conv, gaps)):
    next_gcd = gcd(abs(g), abs(gaps[j+1])) if j+1 < len(gaps) else None
    print(j, p, q, g, next_gcd)
\end{lstlisting}

The computation is illustrative only.  The coprimality statement itself is
proved symbolically in \cref{thm:adjacent} and does not depend on numerical
verification.

\section{Audit checklist for future claimed proofs}

A proposed proof should be rejected or repaired if it contains any of the
following transitions without a new theorem.
\begin{enumerate}[label=(\arabic*)]
\item From $\det(\log a_{ij})=0$ to a rational dependence of rows or columns.
  This is the Four Exponentials step.
\item From real-logarithmic equality to congruence equality at a prime.
  Real and $p$-adic logarithms are not interchangeable.
\item From many integer values to automaticity, recurrence, or an algebraic
  generating function without constructing the required finite-state or
  functional structure.
\item From a termwise $p$-adic divisor to an extra leading determinant
  constant without charging the same multiplier in the archimedean bound.
\item From rational approximants to rationality of $x$ without comparing the
  error to their reduced denominator height.
\item From exhaustive computation below a bound to a proof for all $x$.
\item From a conditional theorem under Four Exponentials, Schanuel, or an
  unproved logarithm-product independence statement to an unconditional
  conclusion.
\end{enumerate}

\begin{thebibliography}{99}

\bibitem{AE1944}
L. Alaoglu and P. Erd\H{o}s,
\newblock \emph{On highly composite and similar numbers},
\newblock Trans. Amer. Math. Soc. \textbf{56} (1944), 448--469.
\newblock \url{https://doi.org/10.1090/S0002-9947-1944-0011087-2}.

\bibitem{Baker1975}
A. Baker,
\newblock \emph{Transcendental Number Theory},
\newblock Cambridge University Press, 1975.

\bibitem{Butler2017}
L. A. Butler,
\newblock \emph{A Diophantine approach to the three and four exponentials conjectures},
\newblock Ramanujan J. \textbf{42} (2017), 199--221.
\newblock \url{https://doi.org/10.1007/s11139-015-9728-2}.

\bibitem{Gelfond}
A. O. Gelfond,
\newblock \emph{Transcendental and Algebraic Numbers},
\newblock Dover, 1960.

\bibitem{HardyWright}
G. H. Hardy and E. M. Wright,
\newblock \emph{An Introduction to the Theory of Numbers}, 6th ed.,
\newblock Oxford University Press, 2008.

\bibitem{Harrison2026}
J. Harrison,
\newblock \emph{Additive relations in irrational powers},
\newblock arXiv:2512.04081v3, July 31, 2026.
\newblock \url{https://arxiv.org/abs/2512.04081}.

\bibitem{KaratarakisWiedijk2026}
M. Karatarakis and F. Wiedijk,
\newblock \emph{A formalization of the Gelfond--Schneider theorem},
\newblock arXiv:2603.24823, March 25, 2026.
\newblock \url{https://arxiv.org/abs/2603.24823}.

\bibitem{Kawashima2026}
M. Kawashima,
\newblock \emph{Linear independence of periods related to polylogarithms},
\newblock arXiv:2605.17756, May 18, 2026.
\newblock \url{https://arxiv.org/abs/2605.17756}.

\bibitem{Schneider}
T. Schneider,
\newblock \emph{Einf\"uhrung in die transzendenten Zahlen},
\newblock Springer, 1957.

\bibitem{Waldschmidt2023}
M. Waldschmidt,
\newblock \emph{The Four Exponentials Problem and Schanuel's Conjecture},
\newblock in \emph{Mathematics Going Forward}, Lecture Notes in Mathematics
2313, Springer, 2023, pp. 579--592.
\newblock \url{https://doi.org/10.1007/978-3-031-12244-6_39}.

\bibitem{AttachedReport}
V. Reshetnikov et al.,
\newblock \emph{The Alaoglu--Erd\H{o}s Integer-Exponent Conjecture: Machine-Checked
Partial Results, Exact Conditional Structure, Determinant Methods, and the Remaining Barrier},
\newblock attached source file
\texttt{Article/alaoglu-erdos-unified-report/alaoglu-erdos-unified-report.tex},
August 21, 2026.

\end{thebibliography}

\end{document}
